\begin{thm}
\label{thm:constant_2coset}
Let $\mathcal{C} = \{C_i\}_{i=0}^n$ be a B-chain with $B = C_0 \supset
C_1 \supset \cdots \supset C_n = \{1\}$. Then $\delta$ is constant on
each double coset $C_j//C_{j+1}\cdot b//C_{j+1}\cdot C_j//C_{j+1}$ for
all $0 \leq i \leq j \leq n-1$ and all $b \in B\backslash C_i$ if and only if,
for all $0 \leq i \leq n-1$,
\[
\delta(b//C_i) = \frac{\delta(b)}{o(C_i)}m(\mathcal{C}, i, b).
\]
\end{thm}
 
\begin{proof}
Fix $i \geq 0$. By Proposition~\ref{prop:subset_corresp}\ref{prop2.5_ii}, $\delta$ is constant on
$C_j//C_{j+1}\cdot b//C_{j+1}\cdot C_j//C_{j+1}$ for all $i \leq j
\leq n-1$ and all $b \in B\backslash C_i$ if and only if $\delta$ is
constant on 
\begin{multline}\label{eq:mod_n-1} 
\left((C_j//C_{n-1})//(C_{j+1}//C_{n-1})\right)\cdot
\left((b//C_{n-1})//(C_{j+1}//C_{n-1})\right)
\\
\cdot
\left((C_j//C_{n-1})//(C_{j+1}//C_{n-1})\right)
%\begin{equation}
%\\
\mbox{ for all }i \leq j \leq n-2 \mbox{ and all }b \in
B\backslash C_i, 
\end{multline}
and
\begin{equation}
\label{eq:just_n-1}
\delta \mbox{ is constant on } C_{n-1}bC_{n-1} \mbox{ for all }b \in
B\backslash C_i.
\end{equation} 
 
Since $B\backslash C_i$ is a union of $C_{n-1}$-$C_{n-1}$ double cosets,
Lemma~\ref{lem:constant} implies that~\eqref{eq:just_n-1} is equivalent
to $\delta(b//C_{n-1}) = \frac{\delta(b)}{o(C_{n-1})}m(\mathcal{C}, n-1,
b)$ for all $b \in B\backslash C_i$. So if $i = n-1$, the theorem is
proved. 

Suppose that $i < n-1$. By induction on $n$, \eqref{eq:mod_n-1}~is
equivalent to 
\[
\delta\left((b//C_{n-1})//(C_i//C_{n-1})\right) = 
\frac{\delta(b//C_{n-1})}{o(C_i//C_{n-1})}m(\mathcal{C}//C_{n-1}, i,
b//C_{n-1})
\]
for all $b \in B\backslash C_i$. So if~\eqref{eq:mod_n-1} and~\eqref{eq:just_n-1} both hold, we have
\[
\begin{aligned}
\delta(b//C_i) &= \delta \left((b//C_{n-1})//(C_i//C_{n-1})\right) =
\frac{\delta(b//C_{n-1})}{o(C_i)/o(C_{n-1})}m(\mathcal{C}//C_{n-1}, i,
b//C_{n-1})
\\
&= \frac{\delta(b)}{o(C_{n-1})o(C_i)/o(C_{n-1})}m(\mathcal{C}, n-1,
b)\cdot m(\mathcal{C}//C_{n-1}, i, b//C_{n-1}).
\end{aligned}
\]
Hence by Lemma~\ref{lem:chain_num_mult}, $\delta(b//C_i) =
\frac{\delta(b)}{o(C_i)}m(\mathcal{C}, i, b)$.

Conversely, suppose that for all $0 \leq i \leq n$ and all $b \in
B\backslash C_i$, $\delta(b//C_i) =
\frac{\delta(b)}{o(C_i)}m(\mathcal{C}, i, b)$. If $i = n-1$, we have
already shown that $\delta$ is constant on $C_{n-1}bC_{n-1}$. Suppose
that $i < n-1$. Then by Lemma~\ref{lem:chain_num_mult}, 
\[
\begin{aligned}
\delta(b//C_i) &= \frac{\delta(b)}{o(C_i)}m(\mathcal{C}, n-1, b) \cdot
m(\mathcal{C}//C_{n-1}, i, b//C_{n-1})
\\
 &= \frac{o(C_{n-1})}{o(C_i)}\delta(b//C_{n-1})m(\mathcal{C}//C_{n-1},
i, b//C_{n-1}).
\end{aligned}
\]
So 
\[
\delta\left((b//C_{n-1})//(C_i//C_{n-1})\right) =
\frac{\delta(b//C_{n-1})}{o(C_i//C_{n-1})}m(\mathcal{C}//C_{n-1}, i,
b//C_{n-1})
\] 
for all $0 \leq i \leq n-2$ and $b \in B\backslash C_i$, whence
$\mathcal{C}//C_{n-1}$ satisfies the same hypothesis as~$\mathcal{C}$.
Then by induction on $n$, $\delta$ is constant on all double cosets
\begin{multline*}
\left((C_j//C_{n-1})//(C_{j+1}//C_{n-1})\right)\cdot
\left((b//C_{n-1})//(C_{j+1}//C_{n-1})\right)
\\
 \cdot
\left((C_j//C_{n-1})//(C_{j+1}//C_{n-1})\right) = C_j//C_{j+1} \cdot
b//C_{j+1}\cdot C_j//C_{j+1},
\end{multline*} 
for all $i \leq j \leq n-2$ and $b \in
B\backslash C_i$. This establishes the converse. 
 \end{proof}
